\section{Planteamiento del problema}
Los modelos modernos de aprendizaje profundo constituyen estructuras matemáticas de elevada complejidad dimensional. 
Al componerse de múltiples capas intermedias con miles de unidades interconectadas, su dinámica interna resulta difícil 
de estructurar e inspeccionar de forma intuitiva, lo que frecuentemente relega el análisis de las representaciones 
intermedias en favor de un enfoque centrado exclusivamente en la correspondencia directa entre entradas y salidas.
\medskip

Esta limitación es particularmente notoria en las redes neuronales convolucionales (CNN), las cuales operan sobre 
volúmenes masivos de datos organizados en tensores tridimensionales de activación. Los métodos y herramientas 
tradicionales de representación visual suelen verse restringidos por las interfaces gráficas convencionales, 
recurriendo a proyecciones bidimensionales (como imágenes en escala de grises o cuadrículas planas promediadas). 
Si se requiere preservar cada canal de manera individual, es necesario exportar cientos de archivos de imagen aislados. 
En ambos casos, se fractura la continuidad espacial tridimensional intrínseca del tensor, imposibilitando el análisis de 
las relaciones geométricas locales y volumétricas que la red construye entre capas adyacentes.
\medskip

Asimismo, el uso de estos formatos de imagen bidimensionales convencionales introduce deficiencias técnicas en la fidelidad 
y gestión de los datos. Formatos gráficos de uso común, como PNG, operan típicamente con una profundidad de $8$ bits por 
canal, truncando drásticamente los valores continuos en punto flotante de $32$ bits propios de las bibliotecas de aprendizaje 
automático (e.g., PyTorch). A esto se añade una redundancia estructural significativa al duplicar valores escalares en canales 
de color independientes (RGB) para representar un único mapa de características, multiplicando el consumo de almacenamiento 
en disco y la sobrecarga en las operaciones de entrada/salida (I/O) al manipular colecciones atomizadas de miles de archivos independientes por inferencia.
\medskip

Actualmente se manifiesta una brecha técnica entre las plataformas de entrenamiento de modelos y los ecosistemas de 
gráficos por computadora: no existe una infraestructura estandarizada que permita serializar y renderizar las estructuras 
tensoriales volumétricas de forma procedural, directa e interactiva dentro de motores gráficos tridimensionales dedicados.
\medskip

Por consiguiente, el problema que aborda este trabajo radica en la degradación de la topología geométrica y la pérdida de 
precisión numérica al visualizar y manipular tensores de activación mediante formatos e interfaces bidimensionales. A esto se 
suma el elevado coste de almacenamiento e I/O que impone la fragmentación en imágenes rasterizadas y la ausencia de un conducto 
(\textit{pipeline}) integrado que conecte los entornos de inferencia con herramientas de computación gráfica. Como consecuencia, el 
desarrollo, análisis estructural y generación de material de comunicación técnica para modelos de visión computacional como 
Tiny YOLOv2 carecen de metodologías volumétricas capaces de representar con fidelidad y eficiencia sus representaciones latentes tridimensionales.